\section{Testing Strategy}
\label{sec:testing}

\subsection{Testing Objectives}

Testing serves multiple purposes in this project:

\begin{enumerate}
  \item \textbf{Correctness}: Verify that the system produces correct outputs.
  
  \item \textbf{Robustness}: Ensure the system handles edge cases and error conditions gracefully.
  
  \item \textbf{Regression Prevention}: Detect unintended behavior changes during development.
  
  \item \textbf{Documentation}: Tests serve as executable specifications of expected behavior.
  
  \item \textbf{Confidence}: Enable developers to refactor and improve code without fear.
\end{enumerate}

\subsection{Testing Pyramid}

The testing strategy follows the testing pyramid:

\begin{verbatim}
           /\               <-- End-to-End Tests (Few, Slow)
          /  \
         /----\
        / Unit \             <-- Unit Tests (Many, Fast)
       /--------\
\end{verbatim}

\textbf{Unit Tests:} Test individual components in isolation (e.g., model wrapper).

\textbf{Integration Tests:} Test interactions between components (e.g., API endpoints with mocked model).

\textbf{End-to-End Tests:} Test the full system with real data and Docker (few, slow, high confidence).

\subsection{Unit Test Framework}

Tests are implemented using \code{pytest}, chosen for:

\begin{enumerate}
  \item Simplicity: Minimal boilerplate.
  
  \item Fixtures: Reusable test setup and teardown.
  
  \item Assertions: Clear, readable assertion syntax.
  
  \item Plugins: Rich ecosystem for mocking, coverage, and reporting.
  
  \item Parametrization: Concise way to test multiple inputs.
\end{enumerate}

\subsection{Test Organization}

Tests are organized as follows:

\begin{verbatim}
  tests/
    __init__.py
    test_api.py          <-- API endpoint tests
    conftest.py          <-- Shared fixtures (optional)
\end{verbatim}

\subsection{Health Endpoint Tests}

\subsubsection{test\_health\_check\_returns\_200}

\begin{verbatim}
  def test_health_check_returns_200(client):
      response = client.get("/health")
      assert response.status_code == 200
\end{verbatim}

\textbf{Purpose:} Verify that the health endpoint is accessible and returns a successful status code.

\textbf{Coverage:} Basic connectivity and endpoint routing.

\subsubsection{test\_health\_check\_response\_format}

\begin{verbatim}
  def test_health_check_response_format(client):
      response = client.get("/health")
      data = response.json()
      assert "status" in data
      assert "model_loaded" in data
      assert isinstance(data["status"], str)
      assert isinstance(data["model_loaded"], bool)
\end{verbatim}

\textbf{Purpose:} Verify the response structure matches the API schema.

\textbf{Coverage:} Schema validation and serialization.

\subsection{Prediction Endpoint Tests}

\subsubsection{Happy Path: test\_predict\_valid\_input}

\begin{verbatim}
  def test_predict_valid_input(client):
      response = client.post(
          "/predict",
          json={"user_id": "196", "movie_id": "242"},
      )
      assert response.status_code == 200
      data = response.json()
      assert "predicted_rating" in data
      assert 1.0 <= data["predicted_rating"] <= 5.0
\end{verbatim}

\textbf{Purpose:} Verify that valid requests produce valid predictions.

\textbf{Key Assertion:} Predictions are within the required range $[1.0, 5.0]$. We do \textit{not} assert 
exact values, as these depend on the random seed and training details.

\subsubsection{test\_predict\_response\_format}

\begin{verbatim}
  def test_predict_response_format(client):
      response = client.post(
          "/predict",
          json={"user_id": "196", "movie_id": "242"},
      )
      assert response.status_code == 200
      data = response.json()
      assert "user_id" in data
      assert "movie_id" in data
      assert "predicted_rating" in data
      assert "model_version" in data
\end{verbatim}

\textbf{Purpose:} Verify all required response fields are present.

\subsubsection{Validation Tests}

\textbf{test\_predict\_missing\_user\_id:}
\begin{verbatim}
  def test_predict_missing_user_id(client):
      response = client.post("/predict", json={"movie_id": "242"})
      assert response.status_code == 422
\end{verbatim}

\textbf{test\_predict\_missing\_movie\_id:}
\begin{verbatim}
  def test_predict_missing_movie_id(client):
      response = client.post("/predict", json={"user_id": "196"})
      assert response.status_code == 422
\end{verbatim}

\textbf{Purpose:} Verify that invalid requests are rejected with 422 (Unprocessable Entity).

\textbf{Coverage:} Input validation and error handling.

\subsubsection{Model Unavailability Test}

\begin{verbatim}
  def test_predict_model_unavailable(client):
      # Simulate model not being loaded
      # This test would require mocking or environment setup
      # Expected: status_code == 503
\end{verbatim}

\textbf{Purpose:} Verify graceful handling when the model is not loaded.

\textbf{Challenge:} In tests, the model is loaded during startup. Comprehensive testing of this case 
requires either mocking or running tests with a deliberately broken model artifact.

\subsection{Running Tests}

Tests are executed using pytest:

\begin{verbatim}
  # Run all tests
  pytest tests/ -v
  
  # Run with coverage report
  pytest tests/ -v --cov=app --cov-report=html
  
  # Run specific test file
  pytest tests/test_api.py -v
  
  # Run specific test
  pytest tests/test_api.py::TestPredictEndpoint::test_predict_valid_input -v
\end{verbatim}

\subsection{Test Fixtures}

The \code{client} fixture provides a TestClient for making requests without spinning up a live server:

\begin{verbatim}
  @pytest.fixture(scope="module")
  def client():
      with TestClient(app) as c:
          yield c
\end{verbatim}

\textbf{Advantages:}
\begin{itemize}
  \item Fast: No network overhead.
  \item Isolated: Each test runs in a fresh application context.
  \item Deterministic: Tests are reproducible and not affected by network conditions.
\end{itemize}

\subsection{Coverage Goals}

Target coverage metrics:

\begin{itemize}
  \item \textbf{API endpoints}: 100\% of critical paths (health, predict).
  
  \item \textbf{Model wrapper}: All methods tested with both valid and error inputs.
  
  \item \textbf{Schemas}: Validation and serialization tested.
  
  \item \textbf{Error handling}: All error paths tested (missing model, invalid input).
\end{itemize}

\subsection{Continuous Integration}

In production, tests are run automatically:

\begin{enumerate}
  \item \textbf{Pre-commit}: Lightweight tests run locally before committing.
  
  \item \textbf{Pull Request}: Full test suite runs on every pull request.
  
  \item \textbf{Merge}: Tests pass before merging to main branch.
  
  \item \textbf{Deployment}: Tests pass before deploying to production.
\end{enumerate}

This workflow prevents regression and ensures code quality.

\subsection{Limitations of Unit Testing}

While comprehensive, unit tests have limitations:

\begin{enumerate}
  \item \textbf{Model behavior}: Tests verify API correctness but not prediction quality.
  
  \item \textbf{Docker isolation}: Tests run on local runtime, not in Docker.
  
  \item \textbf{Dataset changes}: Tests pass with the current dataset but may fail if data changes.
  
  \item \textbf{Latency}: Tests don't measure performance under load.
\end{enumerate}

These gaps are addressed by integration tests, load tests, and production monitoring.
